TEKST jest w rebuttal\_template.tex


DONE - odroznic typy elastycznosci (formalna definicja) od modyfikacji jakie wprowadzamy (evaluation), (pytanie o types), od razu pytanie o elasticty = resistance line 085 - Janek




- eksperyment inferencja bierzemy grid i dodajemy losowo dodajemy dodatkowe z grida duplikaty (pytanie o redundancje) - Adam
DONE - MAE - zbadano jako metoda treningu, my robimy w inferencji - Adam
DONE - - TokenMix - dodac do suplementu rysunek roznocy, dodac cytowanie
DONE - elasticvit jako augmentation - to nie jest nasz cel, my badamy input elasticity
DONE - [1] zacytowac, opisac czym sie roznimy (kodowanie pozycyjne)
DONE - inne kodowania pozycyjne - puscic oba lub wyjasnic dlaczego nie ma sensu (low priority)
DONE - ElasticViT \% - uzywamy 100\% elastic w innych testach, to jest tradeoff, pokazujemy ekstremanlny przypadek
DONE - elasticity = resistance to perturbations
DONE - Sec 3.1, swap order
DONE - Flexivit citation - fix
DONE - line 312 type imagenet
DONE - 479 rephrase



DONE - complexity - jest skomplikowany przy treningu, natomiast w inferencji to tylko kwestia pos encoding, perturbacje sa na wejsciu


DONE - AVE jest przykladem zastosowania, modelu, np. w polaczeniu z RLem, co wykracza poza ten papier. Inne zastosowania, np. medyczne histopato, embodied vision (fuzja wielu skal).
DONE - typy elasticity - wybralismy najczestsze typy (np. drony, histo nie potrzebuja rotacji), natomiast rotaczje sa bardzo ciekawa opcja na future work bo np. jezdzace roboty moga potrzebowac

DONE - - related work do missing data - sprawdzic czy faktycznie pasuje - Janek
[16] - rzeczywiscie robia dropout tylko w treningu
[19] - table 4 b)
[22] - table 3
[29] - metoda pokazuje ile mozna patchy odrzucic, ale w ustrukturyzowany sposob - nie losowy


- central i edge - dopisac zdanie w intro, te metody to jest weryfikacja czy takie samplowanie sprawdzi sie w praktyce bez wchodzenia w np. trenowanie strategii RLowych efektywnego samplowania, co jest out of scope; chcemy sprawdzic samo samplowanie decoupled from learnable sampling policy (i.e. RL).
- transfer learning - dodac w intro, kwestia motywacji


- wyniki nie zaskakujace - tak, chocdzil o sprawdzenie ilosciowe o ile



poprawki
- patch shake, scale, drop wystepuje wszedzie, ale np. rotation tylko w niektorych
